\chapter{Flows and Matchings}\label{ch:flows}

Network flow models anything that moves through a network with capacities: water through pipes, data through links,
workers to jobs. Its power in contests comes from \emph{modelling}. Assignment problems, disjoint paths, choosing a profitable
set of projects with prerequisites, separating two groups at minimum cost, covering a DAG with few chains: all of these reduce
to a maximum flow, a minimum cut or a matching, and then a single library routine solves them. This chapter develops the
theory that makes those reductions trustworthy (residual graphs, the max-flow min-cut theorem, König's and Hall's theorems,
linear programming duality in the form of potentials) and the algorithms behind the templates, from Dinic's blocking flows to
the Hungarian algorithm, Edmonds' blossoms, Stoer--Wagner and Gomory--Hu trees.

\begin{chapterintro}{\ioi}{Chapters~\ref{ch:graphs} and~\ref{ch:shortest} (BFS, Dijkstra, Bellman--Ford).}
  \item understand residual graphs and augmenting paths, and run Dinic's algorithm;
  \item use the max-flow min-cut theorem to model separation and selection problems;
  \item find disjoint paths, maximum bipartite matchings, minimum vertex covers and path covers;
  \item solve assignment problems with the Hungarian algorithm and min-cost flow;
  \item handle lower bounds, general (non-bipartite) matching, global minimum cuts and all-pairs minimum cuts.
\end{chapterintro}

% ======================================================================================================================
\section[Maximum flow]{Maximum flow \hfill \ioi}\label{sec:maxflow}

A \defn{flow network} is a directed graph with a \emph{capacity} $c(e) \ge 0$ on every edge, a source $s$ and a sink $t$. A
\defn{flow} assigns a value $0 \le f(e) \le c(e)$ to every edge such that, at every vertex other than $s$ and $t$, the flow
coming in equals the flow going out (\emph{conservation}). The value of the flow is the net amount leaving $s$, and the
\emph{maximum flow} problem asks for the largest possible value.

\subsection{Residual graphs and augmenting paths}

The key to every flow algorithm is the \defn{residual graph} of a flow. For every edge $u \to v$ with capacity $c$ and flow
$f$, it contains a \emph{forward} edge $u \to v$ with the unused capacity $c - f$, and a \emph{backward} edge $v \to u$ with
capacity $f$: pushing flow backwards along it cancels flow that was sent earlier. A path from $s$ to $t$ in the residual graph
is an \defn{augmenting path}; pushing its bottleneck (smallest residual capacity) along it increases the flow value and keeps
all constraints. Figure~\ref{fig:residual} shows why backward edges are indispensable: a greedy first path can block the
optimum, and only cancelling part of it recovers.

\begin{figure}[htbp]
\centering
\newcommand{\rgpanel}[2]{%
\begin{tikzpicture}[x=0.95cm, y=0.95cm]
  \node[vtxs, hot] (s) at (0, 0) {$s$}; \node[vtxs] (a) at (1.5, 1) {$a$}; \node[vtxs] (b) at (1.5, -1) {$b$};
  \node[vtxs, good] (t) at (3, 0) {$t$};
  #1
  \node[lbl] at (1.5, -1.65) {#2};
\end{tikzpicture}}
\begin{filmstrip}[3]
  \begin{stepbox}{1. a greedy path}
    \rgpanel{\draw[ded] (s) -- (b); \draw[ded] (a) -- (t);
      \draw[hded] (s) -- (a); \draw[hded] (a) -- (b); \draw[hded] (b) -- (t);}{flow 1 along $s, a, b, t$}
  \end{stepbox}
  \begin{stepbox}{2. residual graph}
    \rgpanel{\draw[ded] (a) to[bend right=12] (s); \draw[ded] (t) to[bend left=12] (b);
      \draw[hded] (s) -- (b); \draw[hded, draw=fgverm, dashed] (b) -- (a); \draw[hded] (a) -- (t);}{backward edge $b \to a$}
  \end{stepbox}
  \begin{stepbox}{3. after augmenting}
    \rgpanel{\draw[ded, draw=fgline!60] (a) -- (b);
      \draw[gded] (s) -- (a); \draw[gded] (a) -- (t); \draw[gded] (s) -- (b); \draw[gded] (b) -- (t);}{flow 2: the $a \to b$
      flow is cancelled}
  \end{stepbox}
\end{filmstrip}
\caption{Why residual graphs have backward edges. All capacities are 1. The first augmenting path uses $a \to b$ and blocks
both remaining paths; in the residual graph, the backward edge $b \to a$ (dashed) lets the second path $s, b, a, t$ reroute the
first unit, giving the maximum flow 2.}
\label{fig:residual}
\end{figure}

A flow is maximum exactly when its residual graph has no augmenting path; the next section proves this. The algorithms differ in
\emph{which} augmenting paths they use.

\begin{itemize}
  \item \textbf{Ford--Fulkerson}: any augmenting path, found by DFS. With integer capacities each augmentation adds at least 1,
        so it runs in $\Oh{m F}$ for a maximum flow of value $F$: fine when $F$ is small, as in matching problems.
  \item \textbf{Edmonds--Karp}: always a \emph{shortest} augmenting path (BFS), which bounds the number of augmentations by
        $\Oh{nm}$ regardless of the capacities: $\Oh{nm^2}$.
  \item \textbf{Dinic}: in each phase, BFS computes the distance \emph{level} of every vertex, and a DFS that only follows edges
        from one level to the next pushes a \emph{blocking flow}, saturating all shortest augmenting paths at once
        (Figure~\ref{fig:dinic}). The distance from $s$ to $t$ grows with every phase, so there are at most $n$ phases.
        Dinic runs in $\Oh{n^2 m}$ in general, $\Oh{m\sqrt n}$ on unit-capacity bipartite graphs, and is extremely fast in
        practice; it is the default.
\end{itemize}

\begin{figure}[htbp]
\centering
\begin{tikzpicture}[x=1.35cm, y=1.1cm]
  \foreach \lv in {0,...,3} {
    \fill[fgsky!12] (\lv - 0.35, -1.75) rectangle (\lv + 0.35, 1.75);
    \node[lbl] at (\lv, -2.05) {level \lv};
  }
  \node[vtx, hot] (s) at (0, 0) {$s$}; \node[vtx] (a) at (1, 1) {$a$}; \node[vtx] (b) at (1, -1) {$b$};
  \node[vtx] (c) at (2, 1) {$c$}; \node[vtx] (d) at (2, -1) {$d$}; \node[vtx, good] (t) at (3, 0) {$t$};
  \draw[bded] (s) -- (a); \draw[bded] (s) -- (b); \draw[bded] (a) -- (c); \draw[bded] (b) -- (d);
  \draw[bded] (c) -- (t); \draw[bded] (d) -- (t);
  \draw[ded, draw=fgline, dashed] (a) -- (b);
  \node[notel, anchor=west, text width=5.2cm] at (3.7, 0) {One phase of Dinic: BFS assigns levels, and only edges from level
    $i$ to level $i + 1$ (blue) belong to the level graph. The edge $a \to b$ stays inside a level and is ignored in this
    phase; blocking-flow DFS never looks at it.};
\end{tikzpicture}
\caption{The level graph of one phase of Dinic's algorithm.}
\label{fig:dinic}
\end{figure}

\begin{problem}{Download Speed (CSES)}
Computers are joined by one-way connections with given speeds. Find the maximum speed at which data can be downloaded from
computer 1 to computer $n$.
\solution The maximum flow from 1 to $n$, with speeds as capacities.
\end{problem}

\begin{code}
void solve() {  // Dinic (cpdinic) is pasted above solve()
    int n, m;
    cin >> n >> m;
    Dinic g(n + 1);
    rep(i, m) {
        ll a, b, c;
        cin >> a >> b >> c;
        g.add_edge(a, b, c);  // a one-way connection of capacity c
    }
    print(g.max_flow(1, n));
}
\end{code}

\usetemplate{Max flow (Dinic) + min cut}{Ctrl+Alt+Y, Y}{cpdinic}

% ======================================================================================================================
\section[Minimum cut]{Minimum cut \hfill \ioi}\label{sec:mincut}

An $s$--$t$ \defn[cut]{cut} splits the vertices into a side $S$ containing $s$ and a side $T$ containing $t$; its capacity is the
total capacity of the edges from $S$ to $T$. Every unit of flow must cross from $S$ to $T$ at least once, so the value of any
flow is at most the capacity of any cut.

\begin{fact}[Max-flow min-cut theorem]
The maximum flow value equals the minimum cut capacity. Moreover, after a max-flow algorithm finishes, the vertices reachable
from $s$ in the residual graph form the $S$ side of a minimum cut.
\end{fact}

\emph{Proof.} Let $f$ be a flow with no augmenting path, and let $S$ be the set of vertices reachable from $s$ in its residual
graph; $t \notin S$. Every edge from $S$ to $T$ is saturated (otherwise its head would be reachable), and every edge from $T$ to
$S$ carries no flow (otherwise its backward residual edge would make its tail reachable). So the net flow across the cut is exactly
its capacity, and since the flow value equals the net flow across any cut, the flow value equals this cut's capacity. As no flow
exceeds any cut, both are optimal (Figure~\ref{fig:mincut}). The same argument shows that ``no augmenting path'' implies
``maximum'', which is why the augmenting-path algorithms are correct.

\begin{figure}[htbp]
\centering
\begin{tikzpicture}[x=1.25cm, y=1.1cm]
  \fill[fgsky!15, rounded corners=10pt] (-0.45, -1.7) -- (-0.45, 1.7) -- (2.55, 1.7) -- (2.55, 0.3) -- (4.75, -0.35)
    -- (4.75, -1.7) -- cycle;
  \node[vtx, hot] (s) at (0, 0) {$s$}; \node[vtx] (a) at (2, 1.2) {$a$}; \node[vtx] (b) at (2, -1.2) {$b$};
  \node[vtx] (c) at (4.2, 1.2) {$c$}; \node[vtx] (d) at (4.2, -1.2) {$d$}; \node[vtx, good] (t) at (6.2, 0) {$t$};
  \draw[ded] (s) -- node[wl] {3/3} (a); \draw[ded] (s) -- node[wl] {1/2} (b);
  \draw[ded] (a) -- node[wl] {1/1} (b); \draw[hded] (a) -- node[wl] {2/2} (c);
  \draw[ded] (b) -- node[wl] {2/3} (d); \draw[ded] (c) -- node[wl] {2/3} (t);
  \draw[hded] (d) -- node[wl] {2/2} (t);
  \node[font=\sffamily\scriptsize, text=fgblue] at (0.3, -1.45) {residual-reachable side $S$};
  \node[notel, anchor=west, text width=4.4cm] at (6.9, 0.2) {Labels are flow/capacity; the flow has value 4. The red edges
    $a \to c$ and $d \to t$ leave $S$, both saturated, with total capacity $2 + 2 = 4$: no flow can do better.};
\end{tikzpicture}
\caption{A maximum flow and a minimum cut of the same value. The shaded side $S = \{s, a, b, d\}$ is what is reachable from $s$
in the residual graph.}
\label{fig:mincut}
\end{figure}

\begin{problem}{Police Chase (CSES)}
Close as few two-way streets as possible so that there is no route from crossing 1 to crossing $n$, and print them.
\solution With every street a capacity-1 edge in both directions, the minimum number of streets to close is the minimum cut. After
the max flow, the streets with one end on each side of the residual-reachable set form a minimum cut.
\end{problem}

\begin{code}
void solve() {  // Dinic (cpdinic) is pasted above solve()
    int n, m;
    cin >> n >> m;
    Dinic g(n + 1);
    vector<pii> street(m);
    for (auto& [a, b] : street) {
        cin >> a >> b;
        g.add_edge(a, b, 1, 1);  // a two-way street: capacity 1 each way
    }
    print(g.max_flow(1, n));  // the size of a minimum cut
    // the cut: streets between the source side and the sink side
    for (auto [a, b] : street)
        if (g.min_cut_side(a) != g.min_cut_side(b)) print(a, b);
}
\end{code}

\subsection{Project selection}

Min cuts shine when a problem asks to \emph{choose} a subset under dependencies. There are $n$ projects with profits $p_i$
(possibly negative), and requirements ``choosing $a$ forces choosing $b$''. Build a network with an edge $s \to i$ of capacity
$p_i$ for every profitable project, an edge $i \to t$ of capacity $-p_i$ for every costly one, and an edge $a \to b$ of infinite
capacity for every requirement (Figure~\ref{fig:projsel}). A finite cut with source side $S$ never cuts a requirement edge, so $S$
is a valid selection, and its capacity is exactly the profit given up plus the cost paid:
\[
  \text{best profit} = \sum_{p_i > 0} p_i - \text{min cut} .
\]
This is the \emph{maximum weight closure} problem, and a surprising number of problems are it in disguise: choosing experiments
that need equipment, rooms that need keys, vertices of a graph that must include all successors.

\begin{figure}[htbp]
\centering
\begin{tikzpicture}[x=1.3cm, y=1.0cm]
  \node[vtx, hot] (s) at (0, 0) {$s$};
  \node[vtx] (p1) at (2, 1.2) {1}; \node[vtx] (p2) at (2, -1.2) {2};
  \node[vtx] (p3) at (4, 1.2) {3}; \node[vtx] (p4) at (4, -1.2) {4};
  \node[vtx, good] (t) at (6, 0) {$t$};
  \draw[ded] (s) -- node[wl] {8} (p1); \draw[ded] (s) -- node[wl] {5} (p2);
  \draw[ded] (p3) -- node[wl] {6} (t); \draw[ded] (p4) -- node[wl] {4} (t);
  \draw[ded, line width=1.4pt] (p1) -- node[wl] {$\infty$} (p3); \draw[ded, line width=1.4pt] (p2) -- node[wl] {$\infty$}
    (p3); \draw[ded, line width=1.4pt] (p2) -- node[wl] {$\infty$} (p4);
  \node[notel, anchor=west, text width=5.0cm] at (6.6, 0) {Projects 1, 2 earn 8 and 5; projects 3, 4 cost 6 and 4. Project 1
    needs 3; project 2 needs 3 and 4. Taking everything earns $13 - 10 = 3$; taking only 1 and 3 earns $8 - 6 = 2$. The
    minimum cut, 10, gives $13 - 10 = 3$.};
\end{tikzpicture}
\caption{Project selection as a minimum cut. Infinite edges can never be cut, so the source side is always a closed set of
projects.}
\label{fig:projsel}
\end{figure}

\begin{code}
void solve() {  // Dinic (cpdinic) is pasted above solve()
    int n, m;  // n projects, m requirements "a needs b"
    cin >> n >> m;
    int s = 0, t = n + 1;
    Dinic g(n + 2);
    ll positive = 0;
    rep(i, 1, n + 1) {
        ll p;
        cin >> p;  // profit, possibly negative
        if (p > 0) g.add_edge(s, i, p), positive += p;
        else g.add_edge(i, t, -p);
    }
    rep(j, m) {
        int a, b;
        cin >> a >> b;
        g.add_edge(a, b, INF);  // choosing a forces choosing b
    }
    // best profit = all positive profits - the cheapest way to give up
    // some of them or pay for the needed costs: a minimum cut
    print(positive - g.max_flow(s, t));
}
\end{code}

\begin{recognize}
``Remove the cheapest set of edges or vertices so that $s$ and $t$ are separated'', ``split into two groups where each item has a
cost for each side and some pairs cost extra if separated'', ``choose a subset with profits and prerequisites''. Whenever the
objective is the cost of \emph{separating} things, think of a minimum cut. Costs on vertices become edges by splitting each vertex
into an in-copy and an out-copy joined by an edge of that cost.
\end{recognize}

% ======================================================================================================================
\section[Disjoint paths]{Disjoint paths \hfill \ioi}\label{sec:disjointpaths}

The maximum number of \emph{edge-disjoint} paths from $s$ to $t$ is the maximum flow with all capacities 1: each unit of flow is one
path, and conversely disjoint paths give a flow. Dually, by \defn{Menger's theorem}, it equals the minimum number of edges whose
removal separates $s$ from $t$. For \emph{vertex-disjoint} paths, split every vertex $v$ other than $s$ and $t$ into $v_{\text{in}}$
and $v_{\text{out}}$ joined by an edge of capacity 1, redirect incoming edges to $v_{\text{in}}$ and outgoing ones from
$v_{\text{out}}$ (Figure~\ref{fig:vsplit}); now a vertex can carry at most one unit of flow.

\begin{figure}[htbp]
\centering
\begin{tikzpicture}
  \node[vtx] (u1) at (0, 0.7) {}; \node[vtx] (u2) at (0, -0.7) {}; \node[vtx, cool] (v) at (1.6, 0) {$v$};
  \node[vtx] (w1) at (3.2, 0.7) {}; \node[vtx] (w2) at (3.2, -0.7) {};
  \draw[ded] (u1) -- (v); \draw[ded] (u2) -- (v); \draw[ded] (v) -- (w1); \draw[ded] (v) -- (w2);
  \draw[->, line width=1.2pt, draw=fgblue] (4.0, 0) -- (5.0, 0);
  \begin{scope}[xshift=5.8cm]
    \node[vtx] (x1) at (0, 0.7) {}; \node[vtx] (x2) at (0, -0.7) {};
    \node[vtx, cool] (vi) at (1.5, 0) {\scriptsize in}; \node[vtx, cool] (vo) at (3.0, 0) {\scriptsize out};
    \node[vtx] (y1) at (4.5, 0.7) {}; \node[vtx] (y2) at (4.5, -0.7) {};
    \draw[ded] (x1) -- (vi); \draw[ded] (x2) -- (vi); \draw[hded] (vi) -- node[wl] {1} (vo);
    \draw[ded] (vo) -- (y1); \draw[ded] (vo) -- (y2);
  \end{scope}
\end{tikzpicture}
\caption{Vertex splitting. The edge $v_{\text{in}} \to v_{\text{out}}$ of capacity 1 (or of the vertex's cost) makes a vertex
capacity into an edge capacity.}
\label{fig:vsplit}
\end{figure}

\begin{problem}{Distinct Routes (CSES)}
Find as many routes from room 1 to room $n$ as possible, such that no one-way teleporter is used by two routes, and print them.
\solution The number of routes is the maximum flow with unit capacities. To print them, repeatedly walk from 1 along edges that
carry flow and have not been used yet: by conservation, the walk can always continue until it reaches $n$, and each walk uses up
one unit of flow.
\end{problem}

\begin{code}
void solve() {  // Dinic (cpdinic) is pasted above solve()
    int n, m;
    cin >> n >> m;
    Dinic g(n + 1);
    vvi out(n + 1);  // edge ids leaving each room
    vi to(2 * m);
    rep(i, m) {
        int a, b;
        cin >> a >> b;
        int id = g.add_edge(a, b, 1);
        out[a].push_back(id), to[id] = b;
    }
    ll k = g.max_flow(1, n);
    print(k);
    // each unit of flow is one route: follow edges that carry flow,
    // using every edge once (a used edge is marked by clearing its flow)
    vector<bool> used(2 * m, false);
    rep(r, k) {
        vi route{1};
        int u = 1;
        while (u != n) {
            for (int id : out[u])
                if (!used[id] && g.flow_on(id) > 0) {
                    used[id] = true, u = to[id];
                    break;
                }
            route.push_back(u);
        }
        print(sz(route));
        print(route);
    }
}
\end{code}

% ======================================================================================================================
\section[Bipartite matching]{Bipartite matching \hfill \ioi}\label{sec:matching}

A \defn{matching} is a set of edges no two of which share a vertex. In a bipartite graph with sides $L$ and $R$, a maximum matching
is a maximum flow in the network $s \to L \to R \to t$ with all capacities 1. The flow view translates into a direct algorithm:
an augmenting path is now an \defn{alternating path} that starts at an unmatched left vertex, alternates between non-matching and
matching edges, and ends at an unmatched right vertex; flipping every edge on it increases the matching by one
(Figure~\ref{fig:augmenting}). \defn{Berge's lemma} says that a matching is maximum exactly when no augmenting path exists.

\begin{figure}[htbp]
\centering
\begin{tikzpicture}
  \begin{scope}
    \foreach \i in {0,...,2} { \node[vtx] (l\i) at (0, -\i*1.1) {}; \node[vtx] (r\i) at (2.2, -\i*1.1) {}; }
    \node[vtx, hot] at (l2) {}; \node[vtx, good] at (r0) {};
    \draw[ged] (l0) -- (r1); \draw[ged] (l1) -- (r2);
    \draw[ed] (l0) -- (r0); \draw[ed] (l1) -- (r1); \draw[ed] (l2) -- (r2);
    \draw[line width=4pt, draw=fgorange, opacity=0.35, line join=round] (l2.center) -- (r2.center) -- (l1.center)
      -- (r1.center) -- (l0.center) -- (r0.center);
    \node[lbl, align=center] at (1.1, -3.0) {matching of size 2 and an\\augmenting path (orange)};
  \end{scope}
  \draw[->, line width=1.2pt, draw=fgblue] (3.1, -1.1) -- (4.1, -1.1) node[midway, above, font=\sffamily\scriptsize,
    text=fgblue] {flip};
  \begin{scope}[xshift=5.0cm]
    \foreach \i in {0,...,2} { \node[vtx] (m\i) at (0, -\i*1.1) {}; \node[vtx] (q\i) at (2.2, -\i*1.1) {}; }
    \draw[ed] (m0) -- (q1); \draw[ed] (m1) -- (q2);
    \draw[ged] (m0) -- (q0); \draw[ged] (m1) -- (q1); \draw[ged] (m2) -- (q2);
    \node[lbl] at (1.1, -3.0) {matching of size 3};
  \end{scope}
\end{tikzpicture}
\caption{An augmenting path runs from an unmatched left vertex (red) to an unmatched right vertex (green), alternating between
edges outside and inside the matching (green). Flipping it grows the matching by one.}
\label{fig:augmenting}
\end{figure}

\defn{Kuhn's algorithm} looks for an augmenting path from every left vertex by DFS, in $\Oh{nm}$ total and ten lines of code.
\defn{Hopcroft--Karp} finds a maximal set of shortest augmenting paths per phase, like Dinic, in $\Oh{m\sqrt n}$; use it for
$n, m$ around $10^5$.

\begin{problem}{School Dance (CSES)}
Given which boys and girls are willing to dance together, form as many dance pairs as possible and print them.
\end{problem}

\begin{code}
void solve() {
    int n, m, k;  // n boys, m girls, k possible pairs
    cin >> n >> m >> k;
    vvi adj(n + 1);
    rep(i, k) {
        int a, b;
        cin >> a >> b;
        adj[a].push_back(b);
    }
    // Kuhn: for every boy, look for an augmenting path by DFS
    vi girl_of(m + 1, 0), seen(m + 1, 0);
    int stamp = 0;
    auto augment = [&](auto&& self, int boy) -> bool {
        for (int g : adj[boy]) {
            if (seen[g] == stamp) continue;
            seen[g] = stamp;
            // g is free, or her partner can move to another girl
            if (!girl_of[g] || self(self, girl_of[g])) {
                girl_of[g] = boy;
                return true;
            }
        }
        return false;
    };
    int pairs = 0;
    rep(boy, 1, n + 1) {
        stamp++;
        if (augment(augment, boy)) pairs++;
    }
    print(pairs);
    rep(g, 1, m + 1) if (girl_of[g]) print(girl_of[g], g);
}
\end{code}

\usetemplate{Bipartite matching (Hopcroft-Karp)}{Ctrl+Alt+Y, Shift+H}{cphopcroft}

% ======================================================================================================================
\section[König, Hall and path covers]{König, Hall and path covers \hfill \ioi}\label{sec:konig}

Matchings are more useful for the theorems they come with than for themselves; each of these turns a hard-looking optimisation
into a matching size.

\begin{itemize}
  \item \textbf{König's theorem}: in a bipartite graph, the size of a minimum \emph{vertex cover} (a set of vertices touching every
        edge) equals the size of a maximum matching.
  \item \textbf{Independent sets}: the complement of a vertex cover is an independent set, so a maximum independent set of a
        bipartite graph has $|L| + |R| - \text{matching}$ vertices.
  \item \textbf{Hall's theorem}: $L$ can be matched completely if and only if every subset $X \subseteq L$ has at least $|X|$
        neighbours. The obstruction is always a small neighbourhood.
  \item \textbf{Edge covers}: with no isolated vertices, a minimum edge cover has $n - \text{matching}$ edges.
\end{itemize}

\emph{Constructing the cover (König).} After a maximum matching, let $Z$ be the vertices reachable from unmatched left vertices by
alternating paths (non-matching edges from left to right, matching edges from right to left). Then $(L \setminus Z) \cup (R \cap Z)$
is a vertex cover of exactly the matching's size (Figure~\ref{fig:konig}): every matching edge contributes exactly one endpoint, and
no edge can go from $L \cap Z$ to $R \setminus Z$, or its right end would be reachable.

\begin{figure}[htbp]
\centering
\begin{tikzpicture}
  \foreach \i in {0,...,2} \node[vtx] (l\i) at (0, -\i*1.0) {\scriptsize L\i};
  \foreach \i in {0,...,2} \node[vtx] (r\i) at (2.4, -\i*1.0) {\scriptsize R\i};
  \node[vtx, hot] at (l0) {\scriptsize L0};
  \node[vtx, hot] at (r0) {\scriptsize R0};
  \node[vtx, dim] at (l2) {\scriptsize L2};
  \draw[ged] (l0) -- (r1); \draw[ged] (l1) -- (r0);
  \draw[ed] (l0) -- (r0); \draw[ed] (l0) -- (r2); \draw[ed] (l2) -- (r0);
  \node[notel, anchor=west, text width=6.6cm] at (3.4, -1.0) {Maximum matching L0--R1, L1--R0 (green); L2 is unmatched.
    Alternating paths from L2 reach R0 and, through R0's partner, L1: $Z = \{\text{L2}, \text{R0}, \text{L1}\}$. The cover
    $(L \setminus Z) \cup (R \cap Z) = \{\text{L0}, \text{R0}\}$ (red) touches every edge and has size 2, the matching size.};
\end{tikzpicture}
\caption{König's construction of a minimum vertex cover from a maximum matching.}
\label{fig:konig}
\end{figure}

\usetemplate{Bipartite minimum vertex cover / maximum independent set (Konig)}{Ctrl+Alt+U, V}{cpkonig}

\subsection{Path covers and Dilworth's theorem}

A \defn{path cover} of a DAG is a set of vertex-disjoint paths that together contain every vertex. Make a bipartite graph with a
left copy and a right copy of every vertex and an edge from $u$'s left copy to $v$'s right copy for every DAG edge $u \to v$. A
matching chooses, for some vertices, a successor on their path, with no two vertices choosing the same successor; every matched
edge glues two paths into one. So the minimum number of paths is $n - \text{maximum matching}$ (Figure~\ref{fig:pathcover}).

\begin{figure}[htbp]
\centering
\begin{tikzpicture}
  \begin{scope}
    \node[vtx] (a) at (0, 0) {1}; \node[vtx] (b) at (1.4, 0.8) {2}; \node[vtx] (c) at (1.4, -0.8) {3};
    \node[vtx] (d) at (2.8, 0) {4}; \node[vtx] (e) at (4.2, 0) {5};
    \draw[gded] (a) -- (b); \draw[gded] (b) -- (d); \draw[gded] (d) -- (e);
    \draw[ded] (a) -- (c); \draw[ded] (c) -- (d);
    \node[lbl] at (2.1, -1.5) {two paths: $1, 2, 4, 5$ and $3$};
  \end{scope}
  \begin{scope}[xshift=6.6cm, yshift=1.2cm]
    \foreach \i in {1,...,5} { \node[vtxs] (L\i) at (0, -\i*0.55 + 0.55) {\i}; \node[vtxs] (R\i) at (2.0, -\i*0.55 + 0.55)
      {\i}; }
    \draw[ged] (L1) -- (R2); \draw[ged] (L2) -- (R4); \draw[ged] (L4) -- (R5);
    \draw[ed] (L1) -- (R3); \draw[ed] (L3) -- (R4);
    \node[lbl] at (1.0, -2.75) {as a matching ($5 - 3 = 2$ paths)};
  \end{scope}
\end{tikzpicture}
\caption{A minimum path cover of a DAG (green) and the corresponding matching between ``left'' (predecessor) and ``right''
(successor) copies.}
\label{fig:pathcover}
\end{figure}

\begin{code}
void solve() {  // HopcroftKarp (cphopcroft) is pasted above solve()
    int n, m;
    cin >> n >> m;
    // left copy: v as a predecessor on a path, right copy: v as a successor
    HopcroftKarp hk(n, n);
    rep(i, m) {
        int a, b;
        cin >> a >> b;
        hk.add_edge(a - 1, b - 1);  // the edge a -> b may join two paths
    }
    // every matched edge glues two paths together: n - matching paths
    print(n - hk.max_matching());
}
\end{code}

If paths may share vertices (a \emph{chain} cover of a partial order), first take the transitive closure, then apply the same
formula. \defn{Dilworth's theorem} states that the minimum number of chains covering a partial order equals the size of its
largest \emph{antichain} (a set of pairwise incomparable elements); for example, the minimum number of non-increasing
subsequences that cover a sequence equals the length of its longest strictly increasing subsequence.

% ======================================================================================================================
\section[The assignment problem]{The assignment problem: the Hungarian algorithm \hfill \ioi}\label{sec:hungarian}

In the \defn{assignment problem}, $n$ workers must be matched to $n$ jobs, worker $i$ doing job $j$ costs $a_{ij}$, and the total
cost is to be minimised. It is a minimum-cost perfect matching in a complete bipartite graph, and the \defn{Hungarian algorithm}
solves it in $\Oh{n^3}$ with \emph{potentials}: numbers $u_i$ for rows and $v_j$ for columns with $a_{ij} - u_i - v_j \ge 0$ for
all cells. For any such potentials, every assignment costs at least $\sum u_i + \sum v_j$ (sum the inequalities over the assigned
cells). The algorithm adjusts the potentials, one row at a time, until the assigned cells are all \emph{tight} ($a_{ij} = u_i +
v_j$); then the assignment's cost equals the lower bound, so it is optimal (Figure~\ref{fig:hungarian}). This is linear
programming duality in its simplest form.

\begin{figure}[htbp]
\centering
\begin{tikzpicture}
  \foreach \rw/\vals in {0/{7,2,6,9}, 1/{4,1,8,5}, 2/{6,5,3,7}, 3/{8,9,4,2}} {
    \foreach \cv [count=\ci from 0] in \vals {
      \pgfmathtruncatemacro{\hit}{(\rw == 0 && \ci == 1) || (\rw == 1 && \ci == 0) || (\rw == 2 && \ci == 2) ||
        (\rw == 3 && \ci == 3) ? 1 : 0}
      \ifnum\hit=1 \node[cellc, good, minimum width=8mm, minimum height=7mm] at (\ci*0.8, -\rw*0.7) {\cv};
      \else \node[cellc, minimum width=8mm, minimum height=7mm] at (\ci*0.8, -\rw*0.7) {\cv}; \fi
    }
  }
  \foreach \rw/\uu in {0/5, 1/4, 2/3, 3/2} \node[font=\sffamily\small\bfseries, text=fgblue] at (-0.8, -\rw*0.7) {\uu};
  \foreach \ci/\vv in {0/0, 1/$-3$, 2/0, 3/0} \node[font=\sffamily\small\bfseries, text=fgorange!80!black] at (\ci*0.8, 0.7)
    {\vv};
  \node[lbl] at (-0.8, 1.15) {$u_i$}; \node[lbl] at (1.2, 1.15) {$v_j$};
  \begin{scope}[xshift=5.2cm]
    \foreach \rw/\vals in {0/{2,0,1,4}, 1/{0,0,4,1}, 2/{3,5,0,4}, 3/{6,10,2,0}} {
      \foreach \cv [count=\ci from 0] in \vals {
        \pgfmathtruncatemacro{\hit}{(\rw == 0 && \ci == 1) || (\rw == 1 && \ci == 0) || (\rw == 2 && \ci == 2) ||
          (\rw == 3 && \ci == 3) ? 1 : 0}
        \ifnum\hit=1 \node[cellc, good, minimum width=8mm, minimum height=7mm] at (\ci*0.8, -\rw*0.7) {\cv};
        \else \node[cellc, minimum width=8mm, minimum height=7mm] at (\ci*0.8, -\rw*0.7) {\cv}; \fi
      }
    }
    \node[lbl] at (1.2, 0.7) {reduced costs $a_{ij} - u_i - v_j$};
  \end{scope}
  \node[notel, anchor=west, text width=4.0cm] at (9.0, -1.05) {All reduced costs are $\ge 0$, and the assigned cells (green) are
    0. The assignment costs $2 + 4 + 3 + 2 = 11 = \sum u + \sum v$, the lower bound: it is optimal.};
\end{tikzpicture}
\caption{The Hungarian algorithm's certificate of optimality. Left: costs with row potentials $u$ (blue) and column potentials $v$
(orange); right: the reduced costs.}
\label{fig:hungarian}
\end{figure}

\begin{problem}{Task Assignment (CSES)}
Assign $n \le 200$ employees to $n$ tasks, one each, minimising the total cost, and print the assignment.
\end{problem}

\begin{code}
void solve() {  // hungarian (cphungarian) is pasted above solve()
    int n;
    cin >> n;
    vector<vector<ll>> cost(n, vector<ll>(n));
    for (auto& row : cost) cin >> row;
    auto [total, task] = hungarian(cost);  // row i gets column task[i]
    print(total);
    rep(i, n) print(i + 1, task[i] + 1);
}
\end{code}

\usetemplate{Hungarian algorithm (min-cost assignment)}{Ctrl+Alt+U, Shift+H}{cphungarian}

For maximum-weight assignments, negate the costs; for rectangular problems ($n$ rows, $m \ge n$ columns), the template assigns every
row a distinct column. The $\Oh{n^3}$ bound beats a general min-cost flow on dense assignment problems by a wide margin.

% ======================================================================================================================
\section[Minimum-cost flow]{Minimum-cost flow \hfill \ioi}\label{sec:mcmf}

When every edge also has a cost per unit of flow, the \defn{minimum-cost flow} problem asks for the cheapest flow of a given value
(or the cheapest among maximum flows). The standard method is \emph{successive shortest paths}: repeatedly push flow along the
cheapest $s$--$t$ path of the residual graph, where backward edges have the negated cost, until the required value is reached.
Each augmentation keeps the flow optimal among flows of its value, because a flow is cost-optimal exactly when its residual graph
has no negative cycle. Backward edges have negative costs, so the shortest paths need Bellman--Ford (the template uses SPFA), or
Dijkstra with \emph{potentials}: after each round, add the distances to the potentials, which makes every residual edge
non-negative again (Johnson's reweighting, Section~\ref{sec:bellman}).

\begin{problem}{Parcel Delivery (CSES)}
Send $k$ parcels from city 1 to city $n$; route $i$ carries at most $r_i$ parcels at cost $c_i$ each. Find the minimum total cost,
or report that it is impossible.
\end{problem}

\begin{code}
void solve() {  // MCMF (cpmcmf) is pasted above solve()
    int n, m, k;
    cin >> n >> m >> k;
    MCMF g(n + 1);
    rep(i, m) {
        ll a, b, r, c;
        cin >> a >> b >> r >> c;  // at most r parcels, c per parcel
        g.add_edge(a, b, r, c);
    }
    auto [flow, cost] = g.run(1, n, k);  // the cheapest way to send k
    print(flow < k ? -1 : cost);
}
\end{code}

\usetemplate{Min-cost max-flow}{Ctrl+Alt+Y, Shift+Y}{cpmcmf}

\begin{trick}[Min-cost flow as a modelling language]
Assignment with capacities (a worker can take up to $c$ jobs), $k$ disjoint paths of minimum total length (CSES Distinct Routes
II), transporting goods between warehouses, and ``choose $k$ non-overlapping intervals with maximum weight'' (a chain of time
points with capacity-$k$ edges and interval edges of negative cost) are all min-cost flows. When costs are convex in the amount
(the second unit costs more than the first), model each unit as a separate parallel edge with its own cost.
\end{trick}

% ======================================================================================================================
\section[Flows with lower bounds]{Flows with lower bounds \hfill \ioi}\label{sec:lowerbounds}

Sometimes an edge must carry \emph{at least} $l(e)$ units. Force the lower bounds first: pretend $l(e)$ units flow along every edge,
keep only the remaining capacity $c(e) - l(e)$, and note that each vertex now has an imbalance (more forced flow in than out, or the
reverse). Add a new source $S'$ with edges to every vertex with excess and a new sink $T'$ with edges from every vertex with deficit,
and an edge $t \to s$ of infinite capacity. A feasible flow exists if and only if the maximum flow from $S'$ to $T'$ saturates all
edges out of $S'$; then the flow on $t \to s$ is a feasible $s$--$t$ flow, and one more max flow from $s$ to $t$ (with the helper
edges removed) maximises it. Circulations with demands, and minimum flows (the smallest feasible flow), use the same construction.

\usetemplate{Max flow with lower bounds}{Ctrl+Alt+Y, Shift+G}{cpflowlb}

% ======================================================================================================================
\section[Matching in general graphs]{Matching in general graphs: blossoms \hfill \ioi}\label{sec:blossom}

In a non-bipartite graph, alternating paths still characterise maximum matchings (Berge's lemma holds for every graph), but finding
them is harder: an odd cycle lets an alternating path reach the same vertex with both parities, which confuses a plain search.
\defn{Edmonds' blossom algorithm} handles this. When the search finds an odd alternating cycle, a \emph{blossom}, it contracts the
whole cycle into a single vertex and continues (Figure~\ref{fig:blossom}); an augmenting path in the contracted graph lifts to one
in the original graph, because an odd cycle can be entered at any vertex and traversed on the correct side. The template runs in
$\Oh{n^3}$, comfortable for $n \le 500$.

\begin{figure}[htbp]
\centering
\begin{tikzpicture}
  \begin{scope}
    \node[vtx, hot] (r) at (0, 0) {r}; \node[vtx] (a) at (1.2, 0) {a}; \node[vtx] (b) at (2.4, 0) {b};
    \node[vtx] (c) at (3.2, 0.9) {c}; \node[vtx] (d) at (3.2, -0.9) {d};
    \node[vtx, good] (x) at (4.6, 0.9) {x};
    \draw[ed] (r) -- (a); \draw[ged] (a) -- (b); \draw[ed] (b) -- (c); \draw[ged] (c) -- (d); \draw[ed] (d) -- (b);
    \draw[ed] (c) -- (x);
    \node[lbl, align=center] at (2.3, -1.7) {an odd alternating cycle\\$b, c, d$ (a blossom)};
  \end{scope}
  \draw[->, line width=1.2pt, draw=fgblue] (5.3, 0) -- (6.3, 0) node[midway, above, font=\sffamily\scriptsize, text=fgblue]
    {contract};
  \begin{scope}[xshift=7.0cm]
    \node[vtx, hot] (r2) at (0, 0) {r}; \node[vtx] (a2) at (1.2, 0) {a};
    \node[vtx, cool, minimum size=10mm] (B) at (2.6, 0) {\scriptsize $B$};
    \node[vtx, good] (x2) at (4.0, 0.9) {x};
    \draw[ed] (r2) -- (a2); \draw[ged] (a2) -- (B); \draw[ed] (B) -- (x2);
    \node[lbl, align=center] at (2.0, -1.7) {now $r, a, B, x$ is an augmenting path;\\expanding $B$ routes it via $b, d, c$};
  \end{scope}
\end{tikzpicture}
\caption{A blossom. The search from the unmatched vertex $r$ reaches the odd cycle $b, c, d$; contracting it into $B$ reveals the
augmenting path to the unmatched vertex $x$, which in the original graph goes $r, a, b, d, c, x$.}
\label{fig:blossom}
\end{figure}

\begin{problem}{Maximum matching in a general graph}
Given an undirected graph with $n \le 500$ vertices, find the size of a maximum matching.
\end{problem}

\begin{code}
void solve() {  // blossom_matching (cpblossom) is pasted above solve()
    int n, m;
    cin >> n >> m;
    vvi adj(n);
    rep(i, m) {
        int a, b;
        cin >> a >> b;
        adj[a - 1].push_back(b - 1);
        adj[b - 1].push_back(a - 1);
    }
    vi mate = blossom_matching(adj);  // mate[v] = partner or -1
    int pairs = 0;
    rep(v, n) if (mate[v] > v) pairs++;
    print(pairs);
}
\end{code}

\usetemplate{Maximum matching in a general graph (Edmonds' blossom)}{Ctrl+Alt+U, Shift+G}{cpblossom}

% ======================================================================================================================
\section[Global and all-pairs minimum cuts]{Global and all-pairs minimum cuts \hfill \ioi}\label{sec:globalcut}

\paragraph{Global minimum cut.} The smallest set of edges (by total weight) whose removal disconnects an undirected graph, with no
fixed source or sink, could be found with $n - 1$ max flows from a fixed vertex. The \defn{Stoer--Wagner algorithm} does better
with no flow at all. Each phase grows a set $A$ from any vertex by repeatedly adding the vertex most tightly connected to $A$ (the
largest total edge weight into $A$, a \emph{maximum adjacency order}). If $s$ and $t$ are the last two vertices added, the weight of
$t$'s connection to everything else is a minimum $s$--$t$ cut. Either that cut is the global minimum, or $s$ and $t$ are on the same
side of it, so they can be merged into one vertex. After $n - 1$ phases the answer is the best cut seen, in $\Oh{n^3}$ with a
matrix.

\begin{code}
void solve() {  // stoer_wagner (cpstoerwagner) is pasted above solve()
    int n, m;
    cin >> n >> m;
    vector<vector<ll>> w(n, vector<ll>(n, 0));
    rep(i, m) {
        ll a, b, c;
        cin >> a >> b >> c;
        w[a - 1][b - 1] += c, w[b - 1][a - 1] += c;  // parallel edges add up
    }
    print(stoer_wagner(w));
}
\end{code}

\usetemplate{Global minimum cut (Stoer-Wagner)}{Ctrl+Alt+U, Shift+A}{cpstoerwagner}

\paragraph{All-pairs minimum cuts: the Gomory--Hu tree.} Remarkably, the $\binom{n}{2}$ minimum cuts of an undirected graph take only
$n - 1$ different values, and they are organised by a tree. A \defn{Gomory--Hu tree} is a weighted tree on the same vertices such that
the minimum cut between any $u$ and $v$ in the graph equals the smallest edge weight on the tree path between them
(Figure~\ref{fig:gomoryhu}). Gusfield's algorithm builds it with $n - 1$ max flows: for $s = 1, \ldots, n - 1$, compute the minimum cut
between $s$ and its current tree parent, attach $s$ with that weight, and re-hang every later vertex that lies on $s$'s side of the cut
under $s$. Afterwards every min-cut query is a path-minimum query in a tree.

\begin{figure}[htbp]
\centering
\begin{tikzpicture}
  \begin{scope}
    \node[vtx] (g1) at (0, 0) {1}; \node[vtx] (g2) at (1.6, 1.0) {2}; \node[vtx] (g3) at (1.6, -1.0) {3};
    \node[vtx] (g4) at (3.4, 1.0) {4}; \node[vtx] (g5) at (3.4, -1.0) {5};
    \draw[ed] (g1) -- node[wl] {3} (g2); \draw[ed] (g1) -- node[wl] {2} (g3); \draw[ed] (g2) -- node[wl] {4} (g3);
    \draw[ed] (g2) -- node[wl] {2} (g4); \draw[ed] (g3) -- node[wl] {3} (g5); \draw[ed] (g4) -- node[wl] {4} (g5);
    \draw[ed] (g3) -- node[wl, pos=0.3] {1} (g4);
    \node[lbl] at (1.7, -1.75) {an undirected capacitated graph};
  \end{scope}
  \begin{scope}[xshift=6.0cm]
    \node[vtx] (t1) at (0, 0) {1}; \node[vtx] (t2) at (1.4, 0) {2}; \node[vtx] (t3) at (2.8, 0) {3};
    \node[vtx] (t4) at (4.2, 0) {4}; \node[vtx] (t5) at (5.6, 0) {5};
    \draw[bed] (t1) -- node[wl] {5} (t2); \draw[bed] (t2) -- node[wl] {8} (t3); \draw[bed] (t3) -- node[wl] {6} (t4);
    \draw[bed] (t4) -- node[wl] {7} (t5);
    \node[lbl] at (2.8, -1.0) {its Gomory--Hu tree};
    \node[notel, anchor=north, text width=5.8cm] at (2.8, -1.3) {min cut$(2, 5) = \min(8, 6, 7) = 6$, min cut$(1, x) = 5$ for
      every $x$: vertex 1 has total incident weight 5.};
  \end{scope}
\end{tikzpicture}
\caption{A graph and its Gomory--Hu tree, computed with four max flows. Every one of the ten pairwise minimum cuts is the smallest
weight on a tree path.}
\label{fig:gomoryhu}
\end{figure}

\begin{code}
void solve() {  // gomory_hu (cpgomoryhu) is pasted above solve()
    int n, m, q;
    cin >> n >> m >> q;
    vector<array<ll, 3>> e(m);
    for (auto& [a, b, c] : e) cin >> a >> b >> c, a--, b--;
    auto [par, w] = gomory_hu(n, e);  // a tree on the vertices
    // min cut(u, v) = the lightest edge on the tree path from u to v
    vi depth(n, 0);
    rep(v, 1, n) {  // depth by walking up; n is small in this problem
        for (int x = v; x != 0; x = par[x]) depth[v]++;
    }
    while (q--) {
        int u, v;
        cin >> u >> v;
        u--, v--;
        ll best = LLONG_MAX;
        while (u != v) {  // climb from the deeper end, taking the minimum
            if (depth[u] < depth[v]) swap(u, v);
            chmin(best, w[u]);
            u = par[u];
        }
        print(best);
    }
}
\end{code}

\usetemplate{Gomory-Hu tree (all-pairs minimum cut)}{Ctrl+Alt+U, A}{cpgomoryhu}

\begin{summary}[Chapter summary]
\begin{itemize}
  \item Residual graphs with backward edges make augmenting paths correct; Dinic's level graphs and blocking flows make them fast.
  \item Max flow = min cut; the residual-reachable set is a minimum cut. Project selection is a minimum cut with infinite edges.
  \item Edge-disjoint paths are unit flows; vertex capacities become edges by splitting vertices.
  \item Bipartite matching: augmenting alternating paths (Kuhn, Hopcroft--Karp); König: vertex cover = matching; DAG path cover =
        $n -$ matching; Dilworth for chains and antichains.
  \item Assignment: the Hungarian algorithm with potentials, certified by duality. Min-cost flow: successive shortest paths.
  \item Lower bounds reduce to a feasibility flow; general matching needs blossoms; Stoer--Wagner and Gomory--Hu trees answer
        global and all-pairs minimum cuts.
\end{itemize}
\end{summary}

\begin{practice}
\ppgroup{Flows and cuts}
\begin{pplist}
  \item Download Speed
  \item Police Chase
  \item Distinct Routes
  \item Distinct Routes II
\end{pplist}
\ppgroup{Matchings and assignment}
\begin{pplist}
  \item School Dance
  \item Task Assignment
  \item Parcel Delivery
\end{pplist}
\end{practice}
